\par\bigskip\Needspace{9\baselineskip}
\pdfbookmark[2]{Second Method, [3]}{ceres-second-3}
\begin{center}[3.]\end{center}
\begingroup\small
\noindent\hspace*{\parindent}[The expression under the integral sign is a periodic function of the two arguments {\AstroSymbols ⚳} and {\AstroSymbols ♃}, or of the corresponding mean anomalies \(M\) and \(M'\), and can be expanded in a series in their multiples. This expansion was carried out separately for each term by the method laid out in Theoria interpolationis methodo nova tractata\refstepcounter{footnote}\hyperlink{ceresIII3-note}{\textsuperscript{*)}}\begingroup
\makeatletter\long\def\@makefntext#1{\noindent[*) #1}\makeatother
\footnotetext{\hypertarget{ceresIII3-note}{}Volume III, p. 265. See also the expansions in the treatment of the perturbations of Pallas.]}
\endgroup{} and will be given here as an example for the term
\[
3)\quad \mathrm{const.}\,rr'Q''\sin(3D+\text{\AstroSymbols ⚳}+3\epsilon'-2\alpha-\AscendingNode),\qquad
\text{where }\mathrm{const.}=\frac{\mu a\sin\mathrm{J}}{4\BareRoot{(1-ee)}.\sin i}\mathpunct{\raisebox{.45ex}{,}}
\]
\[
\text{numerically}\qquad\log\mathrm{const.}=2{,}073\ 9541\text{ in arcseconds,}
\]
For this expansion the circumference was divided into 10 parts, and the numerical values of the functions to be expanded were calculated for \(M\) or \(M'=\pm18^{0},\ \pm54^{0},\ \pm90^{0},\ \pm126^{0},\ \pm162^{0}\), with \(M\) and \(M'\) reckoned from aphelion. To use the formulas given below in Article 2 of the first calculation of the general perturbations of Pallas, set \(m=90^{0}-M\) for the moment. If we wish to expand a function \(T\) into the series
\SourceLeaf{410}
\[
\begin{aligned}
T={}&\begin{aligned}[t]&\alpha_0+\alpha_1\cos m+\alpha_2\cos2m+\ldots\\&+\beta_1\sin m+\beta_2\sin2m+\ldots\end{aligned}\\[4pt]
={}&\begin{aligned}[t]&\alpha_0-\beta_1\cos M-\alpha_2\cos2M+\ldots\\&+\alpha_1\sin M-\beta_2\sin2M-\ldots\end{aligned}
\end{aligned}
\]
and are given
\[
\begin{aligned}
\text{for the values }m={}&0,\ 36^{0},\ 72^{0},\ \ldots\ 288^{0},\ 324^{0}\\
T={}&A_0,\ A_1,\ A_2,\ \ldots\ A_8,\ A_9
\end{aligned}
\]
then by the formulas derived in the passage just cited we have
\[
10\alpha_0=A_0+A_2+A_4+A_6+A_8-(A_1+A_3+A_5+A_7+A_9)
\]
\[
\left.\begin{aligned}
5\alpha_n={}&A_0-A_5+(A_1+A_9-A_4-A_6)\cos na\\
&+(A_2+A_8-A_3-A_7)\cos2na\\[3pt]
5\beta_n={}&(A_1-A_9+A_4-A_6)\sin na\\
&+(A_2-A_8+A_3-A_7)\sin2na
\end{aligned}\right\}\quad n\text{ odd}
\]
\[
\left.\begin{aligned}
5\alpha_n={}&A_0+A_5+(A_1+A_9+A_4+A_6)\cos na\\
&+(A_2+A_8+A_3+A_7)\cos2na\\[3pt]
5\beta_n={}&(A_1-A_9-A_4+A_6)\sin na\\
&+(A_2-A_8-A_3+A_7)\sin2na
\end{aligned}\right\}\quad n\text{ even}
\]
\[
a=36^{0}.
\]

The term 3) to be expanded was then decomposed as follows
\[
\begin{aligned}
&rr'Q''\sin(3D+\text{\AstroSymbols ⚳}+3\epsilon'-2\alpha-\AscendingNode)=\\
&\quad r'\cos3\epsilon'\left\{\begin{aligned}&r\cos2\alpha.\,Q''\sin(3D+\text{\AstroSymbols ⚳}-\AscendingNode)\\&-r\sin2\alpha.\,Q''\cos(3D+\text{\AstroSymbols ⚳}-\AscendingNode)\end{aligned}\right\}\\[4pt]
&\quad+r'\sin3\epsilon'\left\{\begin{aligned}&r\cos2\alpha.\,Q''\cos(3D+\text{\AstroSymbols ⚳}-\AscendingNode)\\&+r\sin2\alpha.\,Q''\sin(3D+\text{\AstroSymbols ⚳}-\AscendingNode)\end{aligned}\right\}
\end{aligned}
\]
and for the expressions \(\displaystyle r\genfrac{}{}{0pt}{}{\cos}{\sin}2\alpha.\,Q''\) the values at all pairs
\[
M\text{ or }M'=\pm18^{0},\ \pm54^{0},\ \pm90^{0},\ \pm126^{0},\ \pm162^{0}
\]
were calculated first.]
\par\endgroup
\begin{center}\rule{.30\linewidth}{.4pt}\end{center}
